\section{Related Work}

\subsection{Representation Learning}

Representation learning is a powerful technique for uncovering hidden features in unlabeled data, gaining popularity over time \cite{scholkopf2021toward,hamilton2017representation,le2020contrastive}. Initial methods, like auto-encoders \cite{zhang2019aet, goyal2017nonparametric} and adversarial learning \cite{sarafianos2019adversarial,radford2015unsupervised,donahue2019large} from image samples, have their challenges. They often demand detailed inputs and hefty computational resources. Moreover, their outputs end up in pixel space instead of the more practical feature space, limiting their usefulness for downstream tasks.

A more contemporary approach to representation learning, contrastive learning \cite{tian2020makes,chuang2020debiased}, has demonstrated greater effectiveness. Contrastive learning operates by bringing similar images together while pushing dissimilar ones apart. Supervised contrastive learning \cite{khosla2020supervised} executes in a fully supervised setting, leveraging label information to improve clustering in embedding space and achieve superior performance. The proposed Supervised Contrastive (SupCon) loss outperforms traditional cross-entropy loss on various datasets and demonstrates robustness to natural corruptions and stability to hyperparameter changes. On the other hand, Unsupervised Data Augmentation (UDA) \cite{xie2020unsupervised}  utilizes semi-supervised learning by enhancing the quality of noise through advanced data augmentation methods, significantly boosting model performance on various tasks. An influential algorithm in this domain, SimCLR \cite{chen2020simple}, employs positive and negative image pairs for contrastive learning in a semi-supervised manner. SimCLR has exhibited remarkable performance, achieving SOTA accuracy on the ImageNet \cite{deng2009imagenet} dataset. However, SimCLR comes with significant computational expenses due to its reliance on large batch sizes and the incorporation of positive and negative pairs. These limitations render SimCLR impractical for many applications.

BYOL, SimSiam, Barlow Twin, and Momentum Contrast (MoCo) \cite{he2019momentum} are put forward as more efficient alternatives to SimCLR. A notable departure from SimCLR is eliminating the need for negative pairs during training. Instead, they leverage a pair of augmented views of the same image, streamlining the training process and reducing computational overhead. This design choice positions these RL algorithms as computationally more efficient options than SimCLR. SimCLRv2 \cite{chen2020big} uses a big pretrained model for a second stage data utilization.  In contrast, our approach eliminates the reliance on large models by maintaining a consistent model architecture across all training phases.

\subsection{Pseudo Labeling}

In semi-supervised learning, pseudo-labeling is used on a highly accurate classifier to label and learn from it. This approach has proven effective in various applications like object detection \cite{li2022rethinking,yan2019semi}, image classification \cite{taherkhani2021self,wu2017semi}, and speech recognition \cite{lugosch2022pseudo}. If pseudo labels are inaccurately assigned, classifiers may learn from erroneous labels, diminishing their performance \cite{arazo2020pseudo}.

Many pseudo-labeling methods focus on enhancing the network through consistency loss \cite{zhang2021flexmatch}, moving away from the traditional approach of retraining the model with pseudo labels \cite{iscen2019label}. Consistency-based methods ensure stable model predictions across various augmented views of the same input \cite{rizve2021defense}. However, consistency loss typically updates the network with a fraction of the cross-entropy loss \cite{wang2023kd}, reducing its impact. This diminished emphasis on consistency loss can hinder the network's optimal update, as it does not fully leverage retraining with pseudo labels \cite{xie2020unsupervised}.

MPL addresses the issue of incorrect pseudo-labeling by applying a meta-learning \cite{lee2019meta} approach. In MPL, the teacher uses its softmax predictions on unlabeled data to instruct the student. These predictions, commonly called soft labels, are a well-established concept in the literature on knowledge distillation \cite{cho2019efficacy,phuong2019towards}. This additional adaptation empowers the teacher to generate more accurate PLs for instructing the student network. MPL employs the same bi-level optimization technique from meta-learning \cite{snell2017prototypical, nichol2018reptile, finn2017model} to iteratively update the teacher-student framework. 

Notably, MPL progressively enhances the student network, while our training focuses more on improving the teacher network. This is because the teacher network receives more frequent feedback and updates - first through feedback from the student and then from the meta loss update. Additionally, introducing a controlled noise in either the teacher or student network is crucial, especially when the two networks are similar. This noise injection boosts the learning process's robustness and generalization \cite{wang2023consistent, li2024nicest}.